\renewcommand{\Anno}{2024}

\renewcommand{\NumeroProblema}{2}

\renewcommand{\Testo}{
Ispirato alla pavimentazione di Ljubljana (come in figura (A) sotto), Giulio pensa al seguente problema. Nove quadrati di lato unitario sono collegati tra di loro mediante cerniere che permettono loro di ruotare. Le figure sottostanti, in cui i punti neri sono le cerniere, rappresentano tre possibili configurazioni.

% SSSUP-FIGURA-DA-RICOSTRUIRE: squares.pdf
% Riferimento storico (non importare): fonti-locali/2024/squares.pdf

Per ogni configurazione, chiamiamo \emph{rombo incluso} ciascuno dei 4 rombi bianchi racchiusi dai 9~quadrati, e chiamiamo \emph{cerchio circoscritto} il più piccolo cerchio che contiene tutti i nove quadrati in quella configurazione.

\begin{enumerate}
\renewcommand{\labelenumi}{(\alph{enumi})}

\item Si pone un elastico attorno alla figura (rappresentato dal tratteggio). Si determini il perimetro $P$ del bordo elastico quando l'area del rombo incluso è $\frac 12$.

\item Sia $A$ l'area racchiusa dall'elastico. Si determini il massimo rapporto $\frac AP$, al variare delle possibili configurazioni.

\item Determinare, al variare della configurazione, il massimo raggio del cerchio circoscritto.

\end{enumerate}
}

\renewcommand{\Risposta}{}

\renewcommand{\Aiutino}{}

\renewcommand{\SoluzioneUfficiale}{}

\renewcommand{\SoluzioniAlternative}{}

\renewcommand{\Osservazioni}{}

\renewcommand{\Commenti}{}
